\documentclass[11pt,a4paper]{article}
\usepackage{fontspec,amsmath,amssymb,amsthm,microtype}
\setmainfont{lmroman10-regular.otf}[BoldFont=lmroman10-bold.otf,ItalicFont=lmroman10-italic.otf,BoldItalicFont=lmroman10-bolditalic.otf]
\usepackage{xcolor}
\usepackage[most]{tcolorbox}
\usepackage{mathrsfs}
\providecommand{\lt}{<}
\providecommand{\gt}{>}
\setmonofont{DejaVuSansMono.ttf}[BoldFont=DejaVuSansMono-Bold.ttf,ItalicFont=DejaVuSansMono-Oblique.ttf,BoldItalicFont=DejaVuSansMono-BoldOblique.ttf,Scale=0.86]
\newfontfamily\leanfallback{STIXTwoMath-Regular.otf}[Scale=0.86]
\definecolor{leanrule}{gray}{0.62}
\definecolor{leanback}{gray}{0.962}
\definecolor{leantitle}{gray}{0.22}
\definecolor{leanmuted}{gray}{0.38}
\newtcolorbox{leanbox}[1]{enhanced,breakable,colback=leanback,colframe=leanrule,
  boxrule=0.4pt,arc=0.6mm,left=2.4mm,right=2.4mm,top=1.8mm,bottom=1.8mm,
  before skip=0.8em,after skip=1em,fonttitle=\small\bfseries,coltitle=leantitle,
  colbacktitle=leanback,titlerule=0.3pt,title={#1}}
\newlength{\leancharwidth}
\newenvironment{leancode}{\par\smallskip\begingroup\ttfamily\footnotesize\raggedright
  \settowidth{\leancharwidth}{M}\setlength{\parskip}{0pt}\setlength{\parindent}{0pt}
  \frenchspacing}{\par\endgroup\smallskip}
\newcommand{\leanline}[2]{\hangindent=\dimexpr#1\leancharwidth+4\leancharwidth\relax
  \hangafter=1\noindent\hspace*{#1\leancharwidth}#2\par}
\newcommand{\leansource}[1]{\par{\footnotesize\color{leanmuted}#1}\par}
\newcommand{\leanlabel}[1]{\par\smallskip\noindent{\small\bfseries #1}\par}
\usepackage[margin=25mm]{geometry}
\usepackage[unicode,hidelinks]{hyperref}
\hypersetup{pdfauthor={LangevinLean Project}}
\setlength{\emergencystretch}{3em}
\setlength{\parskip}{0.4em}
\setlength{\parindent}{0pt}
\allowdisplaybreaks[2]
\raggedbottom
\title{\texorpdfstring{Conditional expectation of the gradient discrepancy}{Conditional expectation of the gradient discrepancy}}
\author{LangevinLean Project}
\date{}
\begin{document}\maketitle
Draft mathematical exposition; these notes may contain mistakes.


Let \(Y\) be an observed random variable and let \(D\) be a square-integrable random vector. Under the stated measurability assumptions, the conditional expectation given \(\sigma(Y)\) has a measurable representation \(e(Y)\). Conditional Jensen gives \(\int\|e\|^2d(Y_\#P)\le\mathbb E\|D\|^2\). Applied to \(D_t=\nabla U(Y_t)-\nabla U(X)\), this produces the state-dependent discrepancy in the relative-density weak equation and preserves the Gaussian energy bound.


\section{\texorpdfstring{A measurable representation of a conditional mean}{A measurable representation of a conditional mean}}


Assume \(D\in L^2(P)\) and \(Y\) is measurable. The conditional expectation \(\mathbb E[D\mid\sigma(Y)]\) is strongly measurable with respect to \(\sigma(Y)\). Doob–Dynkin factorization, under the stated assumptions on the spaces, gives a strongly measurable field \(e\) such that \(e(Y)=\mathbb E[D\mid\sigma(Y)]\) almost surely. Conditional Jensen and integration give

\[\int\|e(y)\|^2d(Y_\#P)(y)=\mathbb E\|\mathbb E[D\mid\sigma(Y)]\|^2\le\mathbb E\|D\|^2.\]

The conditional expectation is square-integrable, so the integral on the left is finite. The field is determined only up to a null set under the law of \(Y\), which does not affect these integrals.

\begin{leanbox}{Lean: \texttt{exists\_\allowbreak{}conditional\_\allowbreak{}mean\_\allowbreak{}with\_\allowbreak{}energy}}
\leansource{\texttt{exists\_\allowbreak{}conditional\_\allowbreak{}mean\_\allowbreak{}with\_\allowbreak{}energy} \quad UnadjustedLangevin/\allowbreak{}LSI/\allowbreak{}ConditionalDrift.lean, lines 30–64}
\begin{leancode}
\leanline{0}{theorem\ exists\_conditional\_mean\_with\_energy}
\leanline{4}{\{Ω\ S\ F\ :\ Type*\}\ [mΩ\ :\ MeasurableSpace\ Ω]\ [mS\ :\ MeasurableSpace\ S]}
\leanline{4}{[NormedAddCommGroup\ F]\ [NormedSpace\ ℝ\ F]\ [CompleteSpace\ F]}
\leanline{4}{[MeasurableSpace\ F]\ [BorelSpace\ F]}
\leanline{4}{(P\ :\ Measure\ Ω)\ [IsProbabilityMeasure\ P]}
\leanline{4}{(Y\ :\ Ω\ →\ S)\ (hY\ :\ Measurable\ Y)\ (D\ :\ Ω\ →\ F)\ (hD\ :\ MemLp\ D\ 2\ P)\ :}
\leanline{4}{∃\ e\ :\ S\ →\ F,\ StronglyMeasurable\ e\ ∧}
\leanline{6}{(fun\ ω\ ↦\ e\ (Y\ ω))\ =ᵐ[P]\ P[D\ |\ MeasurableSpace.comap\ Y\ mS]\ ∧}
\leanline{6}{Integrable\ (fun\ y\ ↦\ ‖e\ y‖\ \textasciicircum{}\ 2)\ (P.map\ Y)\ ∧}
\leanline{6}{(∫\ y,\ ‖e\ y‖\ \textasciicircum{}\ 2\ ∂P.map\ Y)\ ≤\ ∫\ ω,\ ‖D\ ω‖\ \textasciicircum{}\ 2\ ∂P}
\end{leancode}
\leanlabel{How the Lean proof follows the argument}
\texttt{exists\_\allowbreak{}conditional\_\allowbreak{}mean\_\allowbreak{}with\_\allowbreak{}energy}: the intermediate facts include \texttt{hcomp}, \texttt{h\allowbreak{}CE}, \texttt{h\allowbreak{}CEint}, \texttt{hmap}, \texttt{h}, in source order; the source steps apply or rewrite \texttt{hfactor.\allowbreak{}symm}, \texttt{integrable\_\allowbreak{}map\_\allowbreak{}measure}, \texttt{h\allowbreak{}Y.\allowbreak{}aemeasurable}, \texttt{Filter.\allowbreak{}Eventually\allowbreak{}Eq.\allowbreak{}of\_\allowbreak{}eq}, \texttt{Real.\allowbreak{}rpow\_\allowbreak{}two}.
\end{leanbox}

\section{\texorpdfstring{The conditional mean on measurable sets}{The conditional mean on measurable sets}}


Suppose \(D\) is integrable and \(e(Y)=\mathbb E[D\mid\sigma(Y)]\) almost surely. For every measurable observation set \(A\), the pushforward integral formula and the defining property of conditional expectation give

\[\int_A e(y)d(Y_\#P)(y)=\int_{Y^{-1}(A)}e(Y)\,dP=\int_{Y^{-1}(A)}D\,dP.\]

The inverse image \(Y^{-1}(A)\) belongs to \(\sigma(Y)\), so the last equality is a valid conditional-expectation identity. All terms are vector-valued integrals under the stated integrability assumptions.

\begin{leanbox}{Lean: \texttt{conditional\_\allowbreak{}mean\_\allowbreak{}setIntegral}}
\leansource{\texttt{conditional\_\allowbreak{}mean\_\allowbreak{}setIntegral} \quad UnadjustedLangevin/\allowbreak{}LSI/\allowbreak{}ConditionalDrift.lean, lines 70–84}
\begin{leancode}
\leanline{0}{theorem\ conditional\_mean\_setIntegral}
\leanline{4}{\{Ω\ S\ F\ :\ Type*\}\ [mΩ\ :\ MeasurableSpace\ Ω]\ [mS\ :\ MeasurableSpace\ S]}
\leanline{4}{[NormedAddCommGroup\ F]\ [NormedSpace\ ℝ\ F]\ [CompleteSpace\ F]}
\leanline{4}{[MeasurableSpace\ F]\ [BorelSpace\ F]}
\leanline{4}{(P\ :\ Measure\ Ω)\ [IsProbabilityMeasure\ P]}
\leanline{4}{(Y\ :\ Ω\ →\ S)\ (hY\ :\ Measurable\ Y)\ (D\ :\ Ω\ →\ F)\ (hD\ :\ Integrable\ D\ P)}
\leanline{4}{(e\ :\ S\ →\ F)\ (he\ :\ StronglyMeasurable\ e)}
\leanline{4}{(hcond\ :\ (fun\ ω\ ↦\ e\ (Y\ ω))\ =ᵐ[P]\ P[D\ |\ MeasurableSpace.comap\ Y\ mS])}
\leanline{4}{(A\ :\ Set\ S)\ (hA\ :\ MeasurableSet\ A)\ :}
\leanline{4}{(∫\ y\ in\ A,\ e\ y\ ∂P.map\ Y)\ =\ ∫\ ω\ in\ Y\ ⁻¹'\ A,\ D\ ω\ ∂P}
\end{leancode}
\leanlabel{How the Lean proof follows the argument}
\texttt{conditional\_\allowbreak{}mean\_\allowbreak{}set\allowbreak{}Integral}: the source steps apply or rewrite \texttt{set\allowbreak{}Integral\_\allowbreak{}map}, \texttt{he.\allowbreak{}aestrongly\allowbreak{}Measurable}, \texttt{h\allowbreak{}Y.\allowbreak{}aemeasurable}. The \texttt{calc} block checks the successive equalities or inequalities in the argument.
\end{leanbox}

\section{\texorpdfstring{Pairing the conditional mean with an observed test}{Pairing the conditional mean with an observed test}}


Let \(\varphi\) be strongly measurable, and assume \(\langle D,\varphi(Y)\rangle\) is integrable. The function \(\varphi(Y)\) is measurable with respect to \(\sigma(Y)\). The conditional expectation rule for a continuous bilinear pairing therefore gives

\[\mathbb E[\langle D,\varphi(Y)\rangle\mid\sigma(Y)]=\langle\mathbb E[D\mid\sigma(Y)],\varphi(Y)\rangle.\]

Integrate this equality, replace the conditional mean by \(e(Y)\), and change variables by the law of \(Y\). The result is

\[\int\langle e(y),\varphi(y)\rangle d(Y_\#P)(y)=\mathbb E\langle D,\varphi(Y)\rangle.\]

The integrability hypothesis is required for this test; it is not inferred for arbitrary unbounded functions.

\begin{leanbox}{Lean: \texttt{conditional\_\allowbreak{}mean\_\allowbreak{}inner\_\allowbreak{}integral}}
\leansource{\texttt{conditional\_\allowbreak{}mean\_\allowbreak{}inner\_\allowbreak{}integral} \quad UnadjustedLangevin/\allowbreak{}LSI/\allowbreak{}ConditionalDrift.lean, lines 90–115}
\begin{leancode}
\leanline{0}{theorem\ conditional\_mean\_inner\_integral}
\leanline{4}{\{Ω\ S\ F\ :\ Type*\}\ [mΩ\ :\ MeasurableSpace\ Ω]\ [mS\ :\ MeasurableSpace\ S]}
\leanline{4}{[NormedAddCommGroup\ F]\ [InnerProductSpace\ ℝ\ F]\ [CompleteSpace\ F]}
\leanline{4}{[MeasurableSpace\ F]\ [BorelSpace\ F]}
\leanline{4}{(P\ :\ Measure\ Ω)\ [IsProbabilityMeasure\ P]}
\leanline{4}{(Y\ :\ Ω\ →\ S)\ (hY\ :\ Measurable\ Y)\ (D\ :\ Ω\ →\ F)\ (hD\ :\ Integrable\ D\ P)}
\leanline{4}{(e\ :\ S\ →\ F)\ (he\ :\ StronglyMeasurable\ e)}
\leanline{4}{(hcond\ :\ (fun\ ω\ ↦\ e\ (Y\ ω))\ =ᵐ[P]\ P[D\ |\ MeasurableSpace.comap\ Y\ mS])}
\leanline{4}{(φ\ :\ S\ →\ F)\ (hφ\ :\ StronglyMeasurable\ φ)}
\leanline{4}{(hprod\ :\ Integrable\ (fun\ ω\ ↦\ inner\ ℝ\ (D\ ω)\ (φ\ (Y\ ω)))\ P)\ :}
\leanline{4}{(∫\ y,\ inner\ ℝ\ (e\ y)\ (φ\ y)\ ∂P.map\ Y)\ =}
\leanline{6}{∫\ ω,\ inner\ ℝ\ (D\ ω)\ (φ\ (Y\ ω))\ ∂P}
\end{leancode}
\leanlabel{How the Lean proof follows the argument}
\texttt{conditional\_\allowbreak{}mean\_\allowbreak{}inner\_\allowbreak{}integral}: the intermediate facts include \texttt{h\allowbreak{}Ypull}, \texttt{hpull}, in source order; the source steps apply or rewrite \texttt{integral\_\allowbreak{}map\_\allowbreak{}of\_\allowbreak{}strongly\allowbreak{}Measurable}, \texttt{he.\allowbreak{}inner}, \texttt{integral\_\allowbreak{}congr\_\allowbreak{}ae}. The \texttt{calc} block checks the successive equalities or inequalities in the argument.
\end{leanbox}

\section{\texorpdfstring{The gradient discrepancy before conditioning}{The gradient discrepancy before conditioning}}


On the product space of the initial state and independent Gaussian variable, put \(Y_t(x,z)=x-t\nabla U(x)+\sqrt{2t}z\) and define

\[D_t(x,z)=\nabla U\!\left(x-t\nabla U(x)+\sqrt{2t}z\right)-\nabla U(x).\]

Thus \(D_t=\nabla U(Y_t)-\nabla U(X)\). The order of subtraction fixes its sign in the weak equation: the conditional field is \(e_t(y)=\nabla U(y)-\mathbb E[\nabla U(X)\mid Y_t=y]\). The definition alone asserts no integrability or energy bound.

\begin{leanbox}{Lean: \texttt{frozenGradientDiscrepancy}}
\leansource{\texttt{frozenGradientDiscrepancy} \quad UnadjustedLangevin/\allowbreak{}LSI/\allowbreak{}ConditionalDrift.lean, lines 122–123}
\begin{leancode}
\leanline{0}{def\ frozenGradientDiscrepancy\ (U\ :\ E\ →\ ℝ)\ (t\ :\ ℝ)\ (p\ :\ E\ ×\ E)\ :\ E\ :=}
\leanline{2}{gradient\ U\ (ULA.malaProposal\ U\ t\ p.1\ p.2)\ -\ gradient\ U\ p.1}
\end{leancode}
\leanlabel{How the Lean proof follows the argument}
\texttt{frozen\allowbreak{}Gradient\allowbreak{}Discrepancy} is a definition with defining expression \texttt{gradient U (ULA.\allowbreak{}mala\allowbreak{}Proposal U t p.\allowbreak{}1 p.\allowbreak{}2) -\allowbreak{} gradient U p.\allowbreak{}1}. The right-hand side specifies the object used in the mathematical statement.
\end{leanbox}


Earlier results used here: \cite[Section~3.34.1]{LangevinLean2026ULASC}.

\section{\texorpdfstring{The sigma-algebra of the interpolated state}{The sigma-algebra of the interpolated state}}


The observation at elapsed time \(t\) is \(Y_t\), and its sigma-algebra is \(\sigma(Y_t)=\{Y_t^{-1}(A):A\text{ measurable}\}\), interpreted as the generated sigma-algebra. A random variable measurable for this sigma-algebra depends only on the value of \(Y_t\). Conditioning the gradient discrepancy on \(\sigma(Y_t)\) therefore yields a field on the state space. If the gradient is Lipschitz, the observation map is measurable and this sigma-algebra is contained in the ambient one.

\begin{leanbox}{Lean: \texttt{frozenDriftSigmaAlgebra}}
\leansource{\texttt{frozenDriftSigmaAlgebra} \quad UnadjustedLangevin/\allowbreak{}LSI/\allowbreak{}ConditionalDrift.lean, lines 126–127}
\begin{leancode}
\leanline{0}{abbrev\ frozenDriftSigmaAlgebra\ (U\ :\ E\ →\ ℝ)\ (t\ :\ ℝ)\ :\ MeasurableSpace\ (E\ ×\ E)\ :=}
\leanline{2}{MeasurableSpace.comap\ (fun\ p\ :\ E\ ×\ E\ ↦\ ULA.malaProposal\ U\ t\ p.1\ p.2)\ inferInstance}
\end{leancode}
\leanlabel{How the Lean proof follows the argument}
\texttt{frozen\allowbreak{}Drift\allowbreak{}Sigma\allowbreak{}Algebra} is a definitional abbreviation with defining expression \texttt{Measurable\allowbreak{}Space.\allowbreak{}comap (fun p : E \ensuremath{\times} E \ensuremath{\mapsto} ULA.\allowbreak{}mala\allowbreak{}Proposal U t p.\allowbreak{}1 p.\allowbreak{}2) infer\allowbreak{}Instance}. The right-hand side specifies the object used in the mathematical statement.
\end{leanbox}


Earlier results used here: \cite[Section~3.34.1]{LangevinLean2026ULASC}.

\section{\texorpdfstring{The conditional discrepancy and its second moment}{The conditional discrepancy and its second moment}}


Assume \(\nabla U\) is \(L\)-Lipschitz and \(G_0=\int\|\nabla U\|^2d\mu<\infty\). The Gaussian displacement estimate proves that \(D_t\) is square-integrable and \(\mathbb E\|D_t\|^2\le L^2(t^2G_0+2tn)\), where \(n=\dim E\). Apply the measurable conditional-mean result to obtain a strongly measurable field \(e_t\) satisfying

\[e_t(Y_t)=\mathbb E[D_t\mid\sigma(Y_t)],\qquad \int\|e_t(y)\|^2d\rho_t(y)\le t^2L^2G_0+2tnL^2,\]\[\rho_t=(Y_t)_\#(\mu\otimes\mathcal N(0,I)).\]

The equality identifies the field for the same Gaussian interpolation used in the algorithm. Its square-integrability and energy bound follow from conditional Jensen and the preceding joint Gaussian estimate.

\begin{leanbox}{Lean: \texttt{exists\_\allowbreak{}frozen\_\allowbreak{}drift\_\allowbreak{}error}}
\leansource{\texttt{exists\_\allowbreak{}frozen\_\allowbreak{}drift\_\allowbreak{}error} \quad UnadjustedLangevin/\allowbreak{}LSI/\allowbreak{}ConditionalDrift.lean, lines 133–163}
\begin{leancode}
\leanline{0}{theorem\ exists\_frozen\_drift\_error}
\leanline{4}{(μ\ :\ Measure\ E)\ [IsProbabilityMeasure\ μ]\ (U\ :\ E\ →\ ℝ)\ (L\ :\ ℝ≥0)}
\leanline{4}{(hLip\ :\ LipschitzWith\ L\ (gradient\ U))}
\leanline{4}{(hgrad\ :\ Integrable\ (fun\ x\ ↦\ ‖gradient\ U\ x‖\ \textasciicircum{}\ 2)\ μ)}
\leanline{4}{(t\ :\ ℝ)\ (ht\ :\ 0\ ≤\ t)\ :}
\leanline{4}{∃\ err\ :\ E\ →\ E,\ StronglyMeasurable\ err\ ∧}
\leanline{6}{(fun\ p\ :\ E\ ×\ E\ ↦\ err\ (ULA.malaProposal\ U\ t\ p.1\ p.2))\ =ᵐ[μ.prod\ (stdGaussian\ E)]}
\leanline{8}{(μ.prod\ (stdGaussian\ E))[frozenGradientDiscrepancy\ U\ t\ |}
\leanline{10}{frozenDriftSigmaAlgebra\ U\ t]\ ∧}
\leanline{6}{Integrable\ (fun\ y\ ↦\ ‖err\ y‖\ \textasciicircum{}\ 2)\ (interpolationLaw\ U\ μ\ t)\ ∧}
\leanline{6}{(∫\ y,\ ‖err\ y‖\ \textasciicircum{}\ 2\ ∂interpolationLaw\ U\ μ\ t)\ ≤}
\leanline{8}{t\ \textasciicircum{}\ 2\ *\ (L\ :\ ℝ)\ \textasciicircum{}\ 2\ *\ (∫\ x,\ ‖gradient\ U\ x‖\ \textasciicircum{}\ 2\ ∂μ)\ +}
\leanline{10}{2\ *\ t\ *\ (Module.finrank\ ℝ\ E\ :\ ℝ)\ *\ (L\ :\ ℝ)\ \textasciicircum{}\ 2}
\end{leancode}
\leanlabel{How the Lean proof follows the argument}
\texttt{exists\_\allowbreak{}frozen\_\allowbreak{}drift\_\allowbreak{}error}: the intermediate facts include \texttt{Y}, \texttt{D}, \texttt{h\allowbreak{}Y}, \texttt{h\allowbreak{}Dmeas}, \texttt{h\allowbreak{}D}, in source order. The \texttt{calc} block checks the successive equalities or inequalities in the argument. The remaining scalar rearrangements are checked by \texttt{ring}.
\end{leanbox}


Earlier results used here: \cite[Section~3.34.1]{LangevinLean2026ULASC}; \cite[Section~3.36.5]{LangevinLean2026ULASC}.

\section{\texorpdfstring{The discrepancy energy relative to the Gibbs measure}{The discrepancy energy relative to the Gibbs measure}}


Suppose a measurable nonnegative density identifies the interpolation law as \(\rho_t=r_t\nu\). The conditional-mean construction gives a square-integrable field under \(\rho_t\). Changing measure proves integrability of the weighted square under \(\nu\), and

\[\int r_t(y)\|e_t(y)\|^2d\nu(y)=\int\|e_t(y)\|^2d\rho_t(y)\le t^2L^2G_0+2tnL^2.\]

The conditional-mean identity is unchanged. This is the form of the energy estimate that appears in entropy production.

\begin{leanbox}{Lean: \texttt{exists\_\allowbreak{}frozen\_\allowbreak{}drift\_\allowbreak{}error\_\allowbreak{}of\_\allowbreak{}density}}
\leansource{\texttt{exists\_\allowbreak{}frozen\_\allowbreak{}drift\_\allowbreak{}error\_\allowbreak{}of\_\allowbreak{}density} \quad UnadjustedLangevin/\allowbreak{}LSI/\allowbreak{}ConditionalDrift.lean, lines 169–194}
\begin{leancode}
\leanline{0}{theorem\ exists\_frozen\_drift\_error\_of\_density}
\leanline{4}{(μ\ ν\ :\ Measure\ E)\ [IsProbabilityMeasure\ μ]\ (U\ :\ E\ →\ ℝ)\ (L\ :\ ℝ≥0)}
\leanline{4}{(hLip\ :\ LipschitzWith\ L\ (gradient\ U))}
\leanline{4}{(hgrad\ :\ Integrable\ (fun\ x\ ↦\ ‖gradient\ U\ x‖\ \textasciicircum{}\ 2)\ μ)}
\leanline{4}{(t\ :\ ℝ)\ (ht\ :\ 0\ ≤\ t)\ (r\ :\ E\ →\ ℝ)}
\leanline{4}{(hrmeas\ :\ Measurable\ r)\ (hr\ :\ ∀\ x,\ 0\ ≤\ r\ x)}
\leanline{4}{(hlaw\ :\ tiltedMeasure\ ν\ r\ =\ interpolationLaw\ U\ μ\ t)\ :}
\leanline{4}{∃\ err\ :\ E\ →\ E,\ StronglyMeasurable\ err\ ∧}
\leanline{6}{(fun\ p\ :\ E\ ×\ E\ ↦\ err\ (ULA.malaProposal\ U\ t\ p.1\ p.2))\ =ᵐ[μ.prod\ (stdGaussian\ E)]}
\leanline{8}{(μ.prod\ (stdGaussian\ E))[frozenGradientDiscrepancy\ U\ t\ |}
\leanline{10}{frozenDriftSigmaAlgebra\ U\ t]\ ∧}
\leanline{6}{Integrable\ (fun\ y\ ↦\ r\ y\ *\ ‖err\ y‖\ \textasciicircum{}\ 2)\ ν\ ∧}
\leanline{6}{(∫\ y,\ r\ y\ *\ ‖err\ y‖\ \textasciicircum{}\ 2\ ∂ν)\ ≤}
\leanline{8}{t\ \textasciicircum{}\ 2\ *\ (L\ :\ ℝ)\ \textasciicircum{}\ 2\ *\ (∫\ x,\ ‖gradient\ U\ x‖\ \textasciicircum{}\ 2\ ∂μ)\ +}
\leanline{10}{2\ *\ t\ *\ (Module.finrank\ ℝ\ E\ :\ ℝ)\ *\ (L\ :\ ℝ)\ \textasciicircum{}\ 2}
\end{leancode}
\leanlabel{How the Lean proof follows the argument}
\texttt{exists\_\allowbreak{}frozen\_\allowbreak{}drift\_\allowbreak{}error\_\allowbreak{}of\_\allowbreak{}density}: the intermediate facts include \texttt{hweighted}, \texttt{hequiv}, \texttt{h}, in source order; the source steps apply or rewrite \texttt{ENNReal.\allowbreak{}to\allowbreak{}Real\_\allowbreak{}of\allowbreak{}Real}, \texttt{smul\_\allowbreak{}eq\_\allowbreak{}mul}, \texttt{integral\_\allowbreak{}tilted\allowbreak{}Measure}.
\end{leanbox}


Earlier results used here: \cite[Section~3.34.1]{LangevinLean2026ULASC}.

\begin{thebibliography}{ULA-LSI}
\bibitem[ULA-SC]{LangevinLean2026ULASC}
LangevinLean Project. \emph{Unadjusted Langevin under Strong Convexity}. LL-01, 2026. Draft report, content revision d50c0f64815b. \url{https://odakyildiz.github.io/LangevinLean/docs/pdf/UnadjustedLangevin/FullProof.pdf}.
\end{thebibliography}

\end{document}
